\subsection{An explicit division-point norm for the harmonic sign}

Let \(E:y^2=x^3-x\), with the order-four automorphism
\(\iota(x,y)=(-x,iy)\). Fix an odd prime \(\ell\), and put
\(\chi=\chi_{-4}(\ell)\).
For \(Q\in E[\ell]\setminus\{0\}\), define
\[
        F(Q)=\frac{x(Q)^2+1}{2y(Q)}.
\]
Both numerator and denominator are nonzero: a zero denominator
would make \(Q\) two-torsion, and the duplication formula shows
that a zero numerator would make \(2Q=(0,0)\).
The same duplication formula gives
\begin{equation}\label{eq:gaussian-norm-function}
       F(Q)^2=x(2Q),\qquad F(\iota Q)=-iF(Q).
\end{equation}
Choose one representative of each \(\langle\iota\rangle\)-orbit
on \(E[\ell]\setminus\{0\}\), and let
\[
       C_\ell=\prod_{Q\ {\rm representative}}F(Q).
\]
Changing the representatives multiplies \(C_\ell\) by a fourth
root of unity.

\begin{proposition}\label{prop:gaussian-division-norm}
One has the exact identity
\[
                         C_\ell^4=\frac{\chi}{\ell}.
\]
For the action of the residue norm-one torus
\[
 T_\ell=
 \{a\in(\mathbf F_\ell[i])^\times:a\overline a=1\}
\]
on the \(\ell\)-torsion points, the multiplier of \(C_\ell\)
is its quadratic character:
\[
        \frac{a(C_\ell)}{C_\ell}
                    =a^{(\ell-\chi)/2}\in\{1,-1\}.
\]
The notation on the left means the monomial action on the factors,
using the second identity in
\eqref{eq:gaussian-norm-function}; it does not presume an
arithmetic realization of \(a\).
\end{proposition}
\begin{proof}
We recall the division-polynomial computation, including its
normalization. On \(y^2=x^3-x\), the standard addition-law
recursions begin with
\[
 \psi_0=0,\quad \psi_1=1,\quad \psi_2=2y,\quad
 \psi_3=3x^4-6x^2-1,\quad
 \psi_4=4y(x^6-5x^4-5x^2+1)
\]
and are
\[
 \psi_{2m+1}=\psi_{m+2}\psi_m^3-\psi_{m-1}\psi_{m+1}^3,
 \qquad
 \psi_{2m}=\frac{\psi_m}{2y}
       \bigl(\psi_{m+2}\psi_{m-1}^2
                       -\psi_{m-2}\psi_{m+1}^2\bigr).
\]
The apparent denominator cancels in the even recursion.
These are the recursions obtained from the chord-and-tangent
addition formulas. In particular \(\psi_\ell\) has leading
coefficient \(\ell\), degree \((\ell^2-1)/2\), and simple roots
at the abscissae of nonzero \(\ell\)-torsion points, modulo sign.

At \((0,0)\), every even division polynomial is zero.
The odd recursion proves by induction that
\[
                \psi_n(0)=(-1)^{(n-1)/2}\qquad(n\text{ odd}).
\]
Indeed for \(n=2m+1\) with \(m=2k\) the surviving term is
\(-\psi_{2k-1}\psi_{2k+1}^3\), of sign \(+1\);
with \(m=2k+1\) it is
\(\psi_{2k+3}\psi_{2k+1}^3\), of sign \(-1\).
Let \(M=(\ell^2-1)/4\), which is even.
Pairing the abscissae under \(\iota\) gives
\[
       \psi_\ell(x)=
          \ell\prod_{Q\ {\rm representative}}
                       (x^2-x(Q)^2).
\]
Evaluating at zero, and using that \(M\) is even, proves
\[
            \prod_{Q\ {\rm representative}}x(Q)^2
                         =\frac{\chi}{\ell}.
\]
Multiplication by two permutes the orbit set. Equation
\eqref{eq:gaussian-norm-function} now gives the claimed fourth
power of \(C_\ell\).

For the second assertion, identify \(E[\ell]\) with the
two-dimensional \(\mathbf F_\ell\)-space \(\mathbf F_\ell[i]\).
The torus \(T_\ell\) is cyclic of order \(n=\ell-\chi\),
and it acts freely on the nonzero vectors. In the nonsplit
case this follows because every nonzero vector is a field unit.
In the split case a nonidentity element is
\((a,a^{-1})\) with \(a\ne1\); neither eigenvalue is one,
so it also fixes no nonzero vector. The number of torus orbits
is
\[
                \frac{\ell^2-1}{\ell-\chi}=\ell+\chi,
                                      \qquad \ell+\chi\equiv2\pmod4.
\]
Choose a generator \(b\) of \(T_\ell\) with \(b^{n/4}=\iota\).
On each regular \(T_\ell\)-orbit, choose representatives
\(v,bv,\ldots,b^{n/4-1}v\) for its \(\langle\iota\rangle\)-orbits.
Multiplication by \(b\) permutes all these factors except the
last, which acquires the multiplier \(-i\).
Over all \(\ell+\chi\) torus orbits the multiplier is therefore
\[
                       (-i)^{\ell+\chi}=-1.
\]
For \(a=b^r\), it is \((-1)^r=a^{n/2}\), as asserted.
Changing the chosen representatives does not affect the result,
because \(T_\ell\) commutes with \(\iota\).
\end{proof}

For an arithmetic element \(\sigma\in H^{\rm cyc}\), its action
on \(E[\ell]\) is in \(T_\ell\).
The first assertion places \(C_\ell\) in
\(\mu_8\ell^{-1/4}\). Since \(\sigma\) fixes all roots of unity,
\[
        \frac{\sigma(C_\ell)}{C_\ell}
                =i^{-\rho_\ell^I(\sigma)}
                =(-1)^{\rho_\ell^I(\sigma)/2}.
\]
Thus the proposition recovers the arithmetic Gaussian square
relation directly from an explicit division-point product.
The evenness of \(\rho_\ell^I\) on the cyclotomic kernel is
used in the last equality.

For a full-GT element, the finite-cover and harmonic-frame
comparison supplies the torus action and hence the monomial
character in the proposition. It does not by itself identify
that action on the product with the Ihara Kummer action on
\(\ell^{-1/4}\). The outstanding comparison can therefore be
stated on this particular norm tensor: compatibility must be
proved for the product of the functions \(F(Q)\), retaining
the fourth-root normalization. The division-polynomial identity
and the finite torus calculation alone do not establish that
compatibility for full GT.
